% ============================================================================
% Expérience 01 — Sélection par stabilité du clustering
% Rapport des résultats sur les 4 datasets
% ============================================================================
% Compilation : pdflatex rapport_stabilite (depuis experiments/01_stabilite/)
% ============================================================================

\documentclass[11pt, a4paper]{article}
\usepackage[utf8]{inputenc}
\usepackage[T1]{fontenc}
\usepackage[french]{babel}
\usepackage{amsmath, amssymb}
\usepackage[margin=2.5cm]{geometry}
\usepackage{booktabs}
\usepackage{float}
\usepackage{graphicx}
\usepackage{xcolor}
\usepackage{hyperref}
\hypersetup{colorlinks=true, linkcolor=blue!60!black}

\newcommand{\R}{\mathbb{R}}
\newcommand{\Dw}{D_w}

\title{Expérience 01 — Sélection par stabilité du clustering}
\author{Pierre Chambet}
\date{Mars 2026}

\begin{document}
\maketitle

\begin{abstract}
On teste la stabilité bootstrap comme critère de sélection du triplet
$(k, \alpha, \omega)$ pour le clustering mixte fonctionnel + vectoriel
(Dw + PAM). Grille fine $21 \times 21$, $B = 150$ répétitions. Sur quatre
jeux de données labellisés, la stabilité sélectionne $k = 2$ pour tous les
datasets. On analyse les triplets retenus, les ARI obtenus, les matrices de
confusion et la sensibilité au choix de $B$.
\end{abstract}

\tableofcontents
\newpage

% ============================================================================
\section{Contexte et objectif}
% ============================================================================

Le rapport de synthèse identifie la \textbf{sélection automatique des
paramètres} $(\alpha, \omega)$ et du nombre de clusters $k$ comme verrou
principal en clustering non supervisé. La silhouette, utilisée par défaut,
conduit souvent à des choix sous-optimaux (ex.\ sur Canadian Weather :
silhouette préfère $\omega = 0$ alors que l'ARI-optimal est proche de
$\omega = 0.8$).

L'idée testée ici : le \textbf{bon triplet} $(k, \alpha, \omega)$ est celui
qui produit des clusters \textbf{reproductibles} lorsqu'on perturbe les
données. On mesure cette reproductibilité par bootstrap.

% ============================================================================
\section{Protocole expérimental}
% ============================================================================

\subsection{Algorithme de stabilité}

Pour chaque triplet $(k, \alpha, \omega)$ dans une grille :

\begin{enumerate}
  \item Calculer la partition $\mathcal{C}$ sur les $n$ individus complets
        (PAM avec la distance $\Dw(\alpha, \omega)$).
  \item Répéter $B$ fois :
  \begin{itemize}
    \item Sous-échantillonner aléatoirement 80\,\% des individus.
    \item Recalculer $\Dw$ et la partition $\mathcal{C}_b^*$ sur ce
          sous-échantillon.
    \item Mesurer l'ARI entre $\mathcal{C}$ (restreinte aux individus
          présents) et $\mathcal{C}_b^*$.
  \end{itemize}
  \item Stabilité $= \frac{1}{B} \sum_{b=1}^{B} \mathrm{ARI}(\mathcal{C},
        \mathcal{C}_b^*)$.
\end{enumerate}

Le triplet avec stabilité maximale est retenu.

\subsection{Paramètres utilisés}

\begin{itemize}
  \item \textbf{Grille} : $\alpha, \omega \in \{0, 0.05, 0.1, \ldots, 1\}$
        (21 valeurs chacun, pas 0.05), $k \in \{2, 3, 4, 5, 6\}$.
  \item \textbf{Bootstrap} : $B = 150$ répétitions, sous-échantillon =
        80\,\%.
  \item \textbf{Datasets} : Canadian Weather, AEMET, Growth, Tecator.
\end{itemize}

Au total : $21 \times 21 \times 5 = 2205$ triplets par dataset, chacun
nécessitant 151 exécutions de PAM (1 partition complète + 150 sur
sous-échantillons).

\subsection{Choix des cas}

Les quatre datasets couvrent des contextes variés :
\begin{itemize}
  \item \textbf{Canadian Weather} (35 stations, 4 régions) : météo,
        forte structure géographique.
  \item \textbf{AEMET} (73 stations, 4 zones) : météo Espagne.
  \item \textbf{Growth} (93 enfants, 2 classes M/F) : courbes de croissance.
  \item \textbf{Tecator} (215 spectres, 3 niveaux de gras) : spectroscopie
        NIR.
\end{itemize}

On les a choisis car (1) ils sont labellisés (évaluation par ARI possible),
(2) ils diffèrent par $n$, $k$ réel et nature des données, (3) le rapport de
synthèse les utilise déjà pour les stratégies A, B, C.

% ============================================================================
\section{Résultats principaux}
% ============================================================================

\subsection{Tableau synthétique}

\begin{table}[H]
\centering
\renewcommand{\arraystretch}{1.2}
\begin{tabular}{lccccccc}
\toprule
\textbf{Dataset} & $n$ & $k_{\mathrm{vrai}}$ & $k_{\mathrm{stab}}$ &
$\alpha$ & $\omega$ & \textbf{Stab.} & \textbf{ARI} \\
\midrule
Canadian & 35 & 4 & 2 & 0.30 & 0.00 & 0.973 & 0.541 \\
AEMET    & 73 & 4 & 2 & 0.00 & 1.00 & 0.989 & 0.267 \\
Growth   & 93 & 2 & 2 & 1.00 & 0.45 & 0.983 & 0.719 \\
Tecator  & 215 & 3 & 2 & 0.40 & 0.30 & 0.985 & 0.467 \\
\bottomrule
\end{tabular}
\caption{Triplet stabilité-optimal et ARI obtenu pour chaque dataset
(grille $21 \times 21$, $B = 150$).}
\label{tab:resumes}
\end{table}

\subsection{Observations immédiates}

\begin{itemize}
  \item Avec $B = 150$ et la grille fine, la stabilité sélectionne
        \textbf{toujours} $k = 2$ sur les quatre datasets.
  \item Les triplets $(\alpha, \omega)$ diffèrent : Canadian favorise
        $(\alpha=0.3, \omega=0)$ (fonctionnel type $D_1$), AEMET choisit
        $(\alpha=0, \omega=1)$ (vectoriel seul), Growth $(\alpha=1, \omega=0.45)$,
        Tecator $(\alpha=0.4, \omega=0.3)$ (mix).
  \item L'ARI varie de 0.27 (AEMET) à 0.72 (Growth). Growth, seul dataset
        avec $k_{\mathrm{vrai}} = 2$, obtient le meilleur ARI.
\end{itemize}

% ============================================================================
\section{Matrices de confusion}
% ============================================================================

Pour chaque dataset, la matrice de confusion croise les \textbf{vraies
classes} (lignes) et les \textbf{clusters prédits} (colonnes) par PAM avec
le triplet stabilité-optimal.

\subsection{Canadian Weather ($k_{\mathrm{stab}} = 2$)}

\begin{table}[H]
\centering
\begin{tabular}{lrr}
\toprule
& Cluster 1 & Cluster 2 \\
\midrule
Arctic     & 0 & 3 \\
Atlantic   & 15 & 0 \\
Continental & 1 & 11 \\
Pacific    & 1 & 4 \\
\midrule
Total      & 17 & 18 \\
\bottomrule
\end{tabular}
\caption{Canadian : 4 régions réduites à 2 clusters. Atlantic (15/15) et
Arctic (3/3) bien séparés ; Continental et Pacific partiellement mélangés.}
\end{table}

\textbf{Interprétation} : Le cluster 1 regroupe Atlantic (15/15) + une partie
de Continental et Pacific. Le cluster 2 regroupe Arctic (3/3) + majorité
de Continental (11), Pacific (4). La fusion de 4 classes en 2 brouille la
structure fine des régions climatiques.

\subsection{AEMET ($k_{\mathrm{stab}} = 2$)}

\begin{table}[H]
\centering
\begin{tabular}{lrr}
\toprule
& Cluster 1 & Cluster 2 \\
\midrule
Atlantique   & 13 & 0 \\
Continental  & 20 & 1 \\
Méditerranée & 6 & 12 \\
Sud         & 3 & 18 \\
\midrule
Total       & 42 & 31 \\
\bottomrule
\end{tabular}
\caption{AEMET : $\omega = 1$ (Ds seul). Les clusters suivent surtout la
géographie/altitude.}
\end{table}

\textbf{Interprétation} : La stabilité choisit $\omega = 1$, i.e.\ la
distance purement vectorielle (altitude, lat, lon). Atlantique et Continental
se retrouvent dans le cluster 1 ; Méditerranée et Sud se partagent entre les
deux. L'ARI faible (0.27) reflète un mélange important.

\subsection{Growth ($k_{\mathrm{stab}} = 2$)}

\begin{table}[H]
\centering
\begin{tabular}{lrr}
\toprule
& Cluster 1 & Cluster 2 \\
\midrule
Garçon & 38 & 1 \\
Fille  & 6  & 48 \\
\midrule
Total  & 44 & 49 \\
\bottomrule
\end{tabular}
\caption{Growth : très bon accord (ARI = 0.72).}
\end{table}

\textbf{Interprétation} : Seul dataset où $k_{\mathrm{stab}} = k_{\mathrm{vrai}}$.
Les clusters correspondent presque aux classes M/F : 38 garçons et 48 filles
bien classés, 7 erreurs (6 filles en cluster 1, 1 garçon en cluster 2).
Le triplet $(\alpha=1, \omega=0.45)$ privilégie la forme des courbes (D1).

\subsection{Tecator ($k_{\mathrm{stab}} = 2$)}

\begin{table}[H]
\centering
\begin{tabular}{lrr}
\toprule
& Cluster 1 & Cluster 2 \\
\midrule
Maigre  & 4  & 108 \\
Moyen   & 32 & 30 \\
Gras    & 41 & 0 \\
\midrule
Total   & 77 & 138 \\
\bottomrule
\end{tabular}
\caption{Tecator : 3 classes fusionnées en 2. Gras vs reste.}
\end{table}

\textbf{Interprétation} : Le cluster 1 isole Gras (41/41) + une partie de
Moyen. Le cluster 2 regroupe Maigre (108/112) et le reste de Moyen. La
frontière Maigre/Moyen est floue ; la stabilité privilégie la séparation
Gras vs non-Gras.

% ============================================================================
\section{Ce qui fonctionne}
% ============================================================================

\begin{itemize}
  \item La stabilité fournit une \textbf{mesure reproductible} : les
        triplets retenus donnent des partitions stables sous bootstrap.
  \item Sur \textbf{Growth}, où $k_{\mathrm{vrai}} = 2$, le triplet
        stabilité-optimal $(k=2, \alpha=1, \omega=0.45)$ obtient un très bon
        ARI (0.72) et une confusion faible.
  \item Les heatmaps de stabilité (grille fine $21 \times 21$) montrent des
        \textbf{zones cohérentes} : pour un $k$ donné, les régions de haute
        stabilité se regroupent. Le symbole $\times$ indique le meilleur
        triplet par $k$.
\end{itemize}

% ============================================================================
\section{Ce qui ne fonctionne pas (ou mal)}
% ============================================================================

\begin{itemize}
  \item \textbf{Canadian, AEMET, Tecator} : $k_{\mathrm{stab}} = 2$ alors que
        $k_{\mathrm{vrai}} \in \{3, 4\}$ ; la fusion brouille la structure.
  \item \textbf{AEMET} : ARI faible (0.27), correspondance partielle avec les
        4 zones réelles.
  \item La stabilité reste \textbf{sensible au choix de $B$} : Canadian varie
        fortement selon $B$ (voir §~\ref{sec:sensibilite-B}).
\end{itemize}

% ============================================================================
\section{Sensibilité au choix de $B$}
\label{sec:sensibilite-B}
% ============================================================================

Des exécutions préalables avec différentes valeurs de $B$ (grilles 11$\times$11
et 21$\times$21) ont montré que le triplet optimal \textbf{varie selon $B$},
notamment pour Canadian Weather :

\begin{table}[H]
\centering
\scriptsize
\begin{tabular}{lccc}
\toprule
\textbf{B} & \textbf{Canadian} $(\alpha, \omega)$ & \textbf{AEMET} $k$ &
\textbf{Tecator} $k$ \\
\midrule
50  & $(0.9, 0)$   & 2 & 2 \\
55  & $(0.3, 0)$   & 2 & 2 \\
70  & $(0.9, 0.9)$ & 3 & 4 \\
150 & $(0.3, 0)$   & 2 & 2 \\
\bottomrule
\end{tabular}
\caption{Triplet optimal pour Canadian ; $k$ retenu pour AEMET et Tecator,
selon $B$. Canadian alterne entre $(0.3, 0)$ et $(0.9, 0.9)$.}
\label{tab:sensibilite-B}
\end{table}

\textbf{Interprétation} : Canadian est le plus sensible : avec B=55 ou B=150
(grille fine), le triplet $(0.3, 0)$ l'emporte ; avec B=70, c'est $(0.9, 0.9)$.
La différence de stabilité entre ces triplets est très faible ($\sim 0.001$),
d'où une instabilité du classement. AEMET et Tecator varient en $k$ (2, 3 ou 4)
selon $B$. Un $B$ plus grand (150) avec une grille fine tend à converger vers
$k=2$ pour tous les datasets.

% ============================================================================
\section{Questions soulevées}
% ============================================================================

\begin{enumerate}
  \item \textbf{Convergence avec $B$} : $B = 150$ avec grille fine donne
        $k = 2$ partout. Faut-il poursuivre l'augmentation de $B$ ou accepter
        ce biais structurel vers les petits $k$ ?
  \item \textbf{Stabilité vs silhouette} : comparaison systématique sur les 4
        datasets pour $(k, \alpha, \omega)$ et ARI.
  \item \textbf{Combinaison de critères} : un vote entre stabilité,
        silhouette et éventuellement gap/prediction strength permettrait-il
        de mieux choisir $(k, \alpha, \omega)$ ?
\end{enumerate}

% ============================================================================
\section{Figures}
% ============================================================================

Les heatmaps de stabilité sont générées par \texttt{analyse\_stabilite.R}
et sauvegardées dans \texttt{figures/}. Chaque figure montre, pour un
dataset et pour $k = 2, \ldots, 6$, la stabilité en fonction de
$(\alpha, \omega)$. Le symbole $\times$ indique le meilleur triplet pour
ce $k$. Rouge = haute stabilité, jaune = faible.

\subsection{Canadian Weather}
\begin{center}
\includegraphics[width=0.85\textwidth]{figures/stabilite_heatmap_canadian.png}
\end{center}

\subsection{AEMET}
\begin{center}
\includegraphics[width=0.85\textwidth]{figures/stabilite_heatmap_aemet.png}
\end{center}

\subsection{Growth}
\begin{center}
\includegraphics[width=0.85\textwidth]{figures/stabilite_heatmap_growth.png}
\end{center}

\subsection{Tecator}
\begin{center}
\includegraphics[width=0.85\textwidth]{figures/stabilite_heatmap_tecator.png}
\end{center}

% ============================================================================
\section{Conclusion provisoire}
% ============================================================================

La stabilité bootstrap (grille $21 \times 21$, $B = 150$) est un critère
\textbf{intéressant} (pas de labels, interprétable, reproductible). Elle
sélectionne $k = 2$ pour tous les datasets. Sur Growth ($k_{\mathrm{vrai}} =
2$), elle performe très bien (ARI 0.72). Sur Canadian, AEMET et Tecator, la
fusion en 2 clusters dégrade l'ARI par rapport à un choix oracle.

L'étude de la \textbf{sensibilité à $B$} (tab.~\ref{tab:sensibilite-B}) montre
que Canadian alterne entre $(0.3, 0)$ et $(0.9, 0.9)$ selon $B$ et la grille.
Un $B$ plus grand réduit la variance de l'estimateur mais ne garantit pas
un choix plus « juste » en termes d'ARI.

Pistes : (1) combiner stabilité avec d'autres critères (vote) ; (2) tester
des corrections du biais vers $k = 2$ ; (3) utiliser les heatmaps pour
identifier des zones robustes de $(\alpha, \omega)$.

\end{document}
